% !TeX program = lualatex
\documentclass[11pt, a4paper]{article}
\usepackage[a4paper, margin=2.5cm]{geometry}
\usepackage{fontspec}
\usepackage[english, bidi=basic, provide=*]{babel}
\babelprovide[import, onchar=ids fonts]{english}
\babelfont{rm}{Noto Sans}

\usepackage{amsmath}
\usepackage{xcolor}

\usepackage{listings}
\usepackage{xcolor}

\lstdefinestyle{cppstyle}{
    language=C++,
    basicstyle=\ttfamily\small,
    keywordstyle=\color{blue},
    commentstyle=\color{green!40!black},
    stringstyle=\color{orange},
    breaklines=true,
    numbers=left,
    numberstyle=\tiny\color{gray},
    showstringspaces=false
}
\lstset{style=cppstyle}

\title{Node Physics Functions}
\author{Solomon}

\definecolor{textorange}{HTML}{FFFFFF}
\definecolor{darkbg}{HTML}{000000}

\pagecolor{darkbg}
\color{textorange}

\begin{document}

\paragraph{Velocity using the previous position of the node:}

\begin{equation*}
\mathbf{v} \approx \frac{\mathbf{x}_t - \mathbf{x}_{t - \Delta t}}{\Delta t}
\end{equation*}

\noindent Here, $\mathbf{x}_t$ is the current position and $\mathbf{x}_{t - \Delta t}$ is the previous position.

\paragraph{Calculate the distance between two nodes:} 

\noindent Suppose the nodes have positions \(\mathbf{x}_A\) and \(\mathbf{x}_B\).
\noindent The vector from A to B is:
\[
\mathbf{d}=\mathbf{x}_B-\mathbf{x}_A
\]

\noindent The distance between them is the magnitude of that vector:
\[
L=\|\mathbf{d}\|
=\sqrt{d_x^2+d_y^2+d_z^2}
\]

\noindent In GLM, it can be calculated with:

\begin{lstlisting}
glm::vec3 delta = nodeB.position - nodeA.position;
float length = glm::length(delta);
\end{lstlisting}



\end{document}